\documentclass{article}

\usepackage[a4paper]{geometry}
\usepackage[english]{babel}
\usepackage{parskip}

\usepackage{amsmath, amssymb}
\usepackage{graphicx}
\usepackage{subcaption}
\usepackage{float}
\usepackage[hidelinks]{hyperref}

\usepackage{listings}
\usepackage{xcolor}

\definecolor{codegreen}{rgb}{0,0.6,0}
\definecolor{codegray}{rgb}{0.5,0.5,0.5}
\definecolor{codepurple}{rgb}{0.58,0,0.82}
\definecolor{backcolour}{rgb}{0.95,0.95,0.92}

\lstdefinestyle{mystyle}{
  backgroundcolor=\color{backcolour},
  commentstyle=\color{codegreen},
  keywordstyle=\color{magenta},
  numberstyle=\tiny\color{codegray},
  stringstyle=\color{codepurple},
  basicstyle=\ttfamily\footnotesize,
  breaklines=true,
  captionpos=b,
  keepspaces=true,
  numbers=left,
  numbersep=5pt,
  showstringspaces=false,
  tabsize=2
}
\lstset{style=mystyle}

\newcommand{\CodeDir}{Code}
\newcommand{\PlotDir}{Plots}
\newcommand{\CalcDir}{Calculations}

\title{NURa Hand-in 2}
\author{Silvan Toet (S1827332)}
\date{Date}

\begin{document}
\maketitle

\textit{Write a short description of your solution in each section. Explain your choices. Explicitly draw conclusions from your results.}

\textit{You can refer to selected lines of code where you think it is necessary. However, you must provide the full code relevant to each (sub)question below it.}

\textit{Don't forget to describe your figures and mention what conclusions you draw from them!}


\section{Satellite galaxies around a massive central}

\subsection{(1a) Determining the normalization}
The spherical volume element \(dV\) can be defined as follows:
We know that the volume can be computed in spherical coordinates using the triple integral
\begin{align}
  \iiint_V n(x) dV & = \int_{\phi=0}^{2\pi} \int_{\theta=0}^{\pi} \int_{r=0}^{r_{max}} n(x) r^2 \sin\theta \, dr\,d\theta\,d\phi                                                \\
                   & = 4 \pi A \langle N_{sat} \rangle \int_{x=0}^{x_{max}=5} x^2 \left(\frac{x}{b}\right)^{a-3} \exp(-\left(\frac{x}{b}\right)^c) dx = \langle N_{sat} \rangle \\
  -> A             & = \frac{1}{4 \pi \int_{x=0}^{x_{max}=5} x^2 \left(\frac{x}{b}\right)^{a-3} \exp(-\left(\frac{x}{b}\right)^c) dx}
\end{align}
Since n(x) where x is the radius relative to virial radius \(x=r/r_{vir}\) is only dependent on r, we can easily solve for \(\theta\) and \(\phi\).
Dividing by \(\langle N_{sat} \rangle\) on both sides and rearranging gives us an expression for \(A\).

Next we use Romberg integration (i.e. Richardson extrapolation applied to integration) to numerically integrate the integral in the expression.
Romberg integration first applies the extended trapezoid rule for N intervals, then 2N, 4N, ... up to $2^m\times$ N intervals, where m is the order of the initial approximations.
Next we use the fact that each successive order's leading error term is exactly a 1/4 times the leading error term of the previous order.
Subtracting the two approximations from each other in such a way as to cancel the leading error term results in m-1 new approximations.
We can do this subtraction trick m-1 times. Since after each subtraction we no longer need the two values used for this, we can store the initial approximations and all subsequent improvements in a vector of length m.
At the end of the algorithm the first element of the vector contains the best estimate of the integral and the best estimate for the error is the absolute difference between the first two elements.
Table~\ref{table_A_constant} shows the value of the normalization constant \(A\) to full precision.


\begin{table}[H]
  \centering
  \caption{Normalization constant of the satellite galaxy profile.}
  \label{table_A_constant}
  \begin{tabular}{|c|}
    \hline
    Normalization constant \(A\)     \\
    \hline
    \input{\CalcDir/satellite_A.txt} \\
    \hline
  \end{tabular}
\end{table}

\subsection{Code for the integrator}

\lstinputlisting[language=Python]{\CodeDir/satellites_integrator_code.txt}

\subsection{(1b) Sampling the satellite distribution}

Describe the sampling method used (e.g. rejection sampling).
Explain how the random number generator was implemented and how its randomness was tested.

\begin{figure}[H]
  \centering
  \includegraphics[width=0.85\textwidth]{\PlotDir/my_solution_1b.png}
  \caption{Comparison between the analytical satellite distribution \(N(x)\) and the histogram of 10 000 sampled satellite galaxies in logarithmic bins.}
\end{figure}

Does your sampled distribution agree with the analytical prediction?

\subsection{Code for sampling}

\lstinputlisting[language=Python]{\CodeDir/satellites_sampling_code.txt}

\subsection{(1c) Selecting and sorting satellites}

Describe the selection of 100 galaxies (equal probability, no duplicates, no rejection).
Explain the sorting algorithm you used.

\begin{figure}[H]
  \centering
  \includegraphics[width=0.8\textwidth]{\PlotDir/my_solution_1c.png}
  \caption{Cumulative number of selected galaxies as a function of relative radius.}
\end{figure}

\subsection{Code for selection and sorting}

\lstinputlisting[language=Python]{\CodeDir/satellites_selection_code.txt}

\subsection{(1d) Numerical derivative}

\begin{table}[H]
  \centering
  \caption{Derivative of the satellite number density at \(x=1\).}
  \begin{tabular}{|c|c|}
    \hline
    Method     & Value                                         \\
    \hline
    Analytical & \input{\CalcDir/satellite_deriv_analytic.txt} \\
    Numerical  & \input{\CalcDir/satellite_deriv_numeric.txt}  \\
    \hline
  \end{tabular}
\end{table}

Do your analytical and numerical derivatives agree? Why?

\subsection{Code for derivative calculation}

\lstinputlisting[language=Python]{\CodeDir/satellites_derivative_code.txt}

\section{Heating and cooling in HII regions}


\subsection{(2a) Photoionization heating and radiative recombination cooling}

Describe the implementation of your root-finding method (e.g. bisection, Newton-Raphson).

\begin{table}[H]
  \centering
  \caption{Equilibrium temperature for the simplified heating–cooling model.}
  \begin{tabular}{|c|c|c|}
    \hline
    Temperature [K] & Absolute error & Relative error \\
    \hline
    \input{\CalcDir/equilibrium_temp_simple.txt}
    \\
    \hline
  \end{tabular}
\end{table}

How many iterations it took? Discuss the convergence of your root finder.

\subsection{(2b) More realistic model}


\begin{table}[H]
  \centering
  \caption{Equilibrium temperature for different gas densities.}
  \begin{tabular}{|c|c|}
    \hline
    Density \(n_H\) [cm\(^{-3}\)] & Equilibrium temperature [K]                   \\
    \hline
    $10^{-4}$                     & \input{\CalcDir/equilibrium_low_density.txt}  \\
    $1$                           & \input{\CalcDir/equilibrium_mid_density.txt}  \\
    $10^4$                        & \input{\CalcDir/equilibrium_high_density.txt} \\
    \hline
  \end{tabular}
\end{table}

Discuss the results.  What is the difference for the low, intermediate and high gas density?


\subsection{Code for Exercise 2}

\lstinputlisting[language=Python]{\CodeDir/heatingcooling_code.txt}

\section{Conclusions}

\end{document}
